\documentclass[a4paper, 12pt]{book} % changed from scrbook
% not using \documentclass[a4paper, twoside=semi]{scrreprt} anymore

\usepackage{../Style-Files/InfoSec}
\usepackage{algorithmic}


\definecolor{bg}{RGB}{20,20,20}
\definecolor{fg}{RGB}{220,220,220}
\definecolor{keyword}{RGB}{130,170,255}
\definecolor{comment}{RGB}{120,120,120}
\definecolor{string}{RGB}{255,200,120}

\lstdefinestyle{darkbash}{
    backgroundcolor=\color{bg},
    basicstyle=\scriptsize\ttfamily\color{fg},
    keywordstyle=\color{keyword},
    commentstyle=\color{comment},
    autogobble=true,
    stringstyle=\color{string},
    breaklines=true,
    showstringspaces=false,
    frame=single
}

% ------------------------------------------------
% Dracula background
% ------------------------------------------------
\definecolor{draculabg}{HTML}{282A36}

% ------------------------------------------------
% Minted / Pygments
% ------------------------------------------------
\setminted{
    style=dracula,
    bgcolor=draculabg,
    fontsize=\small,
    fontfamily=tt,
    breaklines=true,
    autogobble=false,
    tabsize=4,
    linenos=false,
    framesep=0pt,
    baselinestretch=1.05
}

% \usepackage{titlesec}
% \setkomafont{chapter}{\Huge\bfseries}
% \setkomafont{section}{\Large\bfseries}
% \setkomafont{subsection}{\large\bfseries}
\usepackage{fancyhdr, lastpage}
\usepackage{lstautogobble}
\usepackage{euscript}


\makeatletter
\providecommand{\subtitle}[1]{\gdef\@subtitle{#1}}
\providecommand{\@subtitle}{}
\makeatother


\usepackage{booktabs}
\usepackage{array}
\usepackage{multirow}




\usepackage{newfloat}
%\usepackage{tocloft}

% % Define new float environment
% \DeclareFloatingEnvironment[
%   fileext=lof,
%   listname={List of Flowcharts},
%   name=Flowchart,
% ]{flowchart}


\pagestyle{fancy}
%\fancyhf{}
\lhead{}
\chead{}
%\rhead{Introduction to the Theory of Computation}
\cfoot{\thepage\ of \pageref{LastPage}}
\rfoot{Version 8.3}




\RequirePackage{etex} 


%--- The documentation for the title can be found at: https://en.wikibooks.org/wiki/LaTeX/Title_Creation
\title{Information Security \\ Mathematical Preliminaries}
\author{Ahmed U. Ashraf}
\date{\today }
% \subtitle{Version~\version}

\begin{document}

% \include{glossary}


\maketitle
%--- we would like to hyperlink each item in contensts page
% \let\hypercontentsline=\contentsline
% \renewcommand{\contentsline}[4]{\hypertarget{toc.#4}{}\hypercontentsline{#1}{#2}{#3}{#4}}
%--- We print ToC: Table of Contents

\DeclareFontShape{U}{stmry}{b}{n}{<-> ssub * stmry/m/n}{}

\tableofcontents\label{toc}
%--- We print the list of figures, this might go at the end.




% \listoffigures
% %--- We print the list of tables, this might also go at the end.
% \listoftables
% %--- This is to add list of symbols as glossary
% % \printunsrtglossary[type=symbols,style=long,title={List of Symbols}]


% % List of Algorithms (- taken the idea from Discrete Algorithmic Mathematics by Stephen B.Maurer and ANthony Ralston.pdf)

% \listofalgorithms
% \addcontentsline{toc}{chapter}{List of Algorithms}

% \listofflowcharts
% %\addcontentsline{toc}{chapter}{List of Flowcharts}


% GENERAL REFERENCES
% Algebra for Applications by Arkadii Slinko


\chapter{Integers}

References are \cite{slinko2015algebra, childs2019cryptology,silverman2008introduction, shemanske2017modern}


Let $a, b \in \mathbb{Z}$ with $a \neq 0$. Then $a$ is called a \emph{divisor}\index{divisor}\footnote{So the term divisor is not defined for $0$.} or \emph{factor}\index{factor} of $b$ if there exists an integer $q \in \mathbb{Z}$ such that $b = qa$, and $b$ is then said to be \emph{divisible}\index{divisible} by $a$, or a \emph{multiple}\index{multiple} of $a$. We denote this by $a|b$. Thus, by definition,
\begin{align}
    a|b ~\overset{\text{def}}{\Leftrightarrow}~\exists~q \in \mathbb{Z} \st b = qa
\end{align}
The integer $q$ is called the \emph{quotient}\index{quotient} in this setting.

\begin{theorem}
    For $a, b, c \in \mathbb{Z}$ with $a \neq 0$,
    \begin{enumerate}
        \item $a|a$, and $1|a$.
        \item $a|b~\Rightarrow~a|-\hspace{-1mm}b$.
        \item $a|b$ and $a|c$~$\Rightarrow~a|(b+c)$.
        \item $a|b$ and $b|c$ ~$\Rightarrow~a|c$.
        \item $a|b$ and $b|a$~$\Rightarrow~a = \pm b$.
    \end{enumerate}
\end{theorem}

Any integer $a$ is trivially divisible by ˙$1$ and by $a$. If an integer $p > 1$  has only these trivial divisors it is called \emph{prime}\index{prime}\footnote{Note that we don't allow \emph{negative} primes.}, otherwise it is called \emph{composite}\index{composite}. The first  primes are 
\begin{align}
    2, 3, 5, 7, 11, 13, \ldots
\end{align}
The number 1 is (by definition) not considered prime. The number 2 is the only even prime\footnote{This is in itself not a profound statement: one might as well make the observation that 8923 has the remarkable property of being the only prime which is divisible by 8923. Nevertheless, the prime 2 seems, in many cases, to call for treatment different from that of odd primes.}.

\begin{theorem}
    Given any integer $a, b$ with $a > 0$, there exists unique integers $q, r$ such that 
    \begin{align}
        b = qa + r
    \end{align}
    with $0 \leq r < a$. In this case, we also say that $q$ and $r$ are, respectively, the quotient and the remainder of $b$ when it is divided by $a$. It is often said that $q$ and $r$ are the quotient and remainder of dividing $a$ into $b$. The notation $r \equiv b \pmod a$ is often used.
\end{theorem}


\begin{definition}
    A positive integer $d$ is called the \emph{greatest common divisor}\index{greatest common divisor} of the integers $a$ and $b$, where $a$ and $b$, not both 0, written $d = \gcd(a, b)$ if the following properties hold:
    \begin{enumerate}
        \item $d$ divides $a$ and $b$, and
        \item if an integer $c$ divides $a$ and $b$, then $c$ must divide $d$.
    \end{enumerate}
    If $gcd(a, b) = 1$, we say that $a$ and $b$ are \emph{relatively prime}\index{relatively prime}.
\end{definition}


\begin{theorem}
    (\emph{Bezout's identity}\index{Bezout's identity}) If $a, b \in \mathbb{Z}$, then
    \begin{align}
        \gcd(a, b) = na + mb
    \end{align} 
    for some integers $a, b \in \mathbb{Z}$. Moreover, $\gcd(a, b)$ is the smallest positive integer which is an integer linear combination $na + mb$ of $a$ and $b$.
\end{theorem}


\section{Fundamental Theorem of Arithmetic}


\begin{definition}
    An integer $p$ is called a prime if $p > 1$ and if the only positive integers dividing $p$ are $1$ and $p$.
\end{definition}


\begin{lemma}
    (\emph{Euclid's lemma}\index{Euclid's lemma}) 
    Let $p$ be a prime number, and suppose that $p$ divides the product $ab$ of two integers $a$ and $b$. Then $p$ divides at least one of $a$ and $b$. More generally, if $p$ divides a product of integers, say
    \begin{align}
        p|a_1 a_2 \cdots a_\ell
    \end{align}
    then $p$ divides at least one of the individual $a_i$.
\end{lemma}
\begin{theorem}
    (\emph{Fundamental Theorem of Arithmetic}\index{Fundamental Theorem of Arithmetic}) Every positive integer $n > 1$ can be expressed as a product of primes (with perhaps only one factor), that is
    \begin{align}
        n = p_1^{k_1} p_2^{k_2} \cdots p_{\ell}^{k_\ell}
    \end{align}
    where $p_1, p_2, \ldots, p_\ell$ are distinct primes and $k_1, k_2, \ldots, k_\ell$ are positive integers. This factorisation is unique apart from the order of the prime factors.
\end{theorem}

\begin{theorem}
    There are infinitely many primes.
\end{theorem}
\begin{proof}
    Assume that there are only a finite number of primes 
    \begin{align}
        p_1, p_2, \ldots, p_\ell
    \end{align}
    Let
    \begin{align}
        n = p_1 p_2 \cdots p_\ell + 1
    \end{align}
    Then $p_i$ does not divide $n$, for all $1 \leq i \leq \ell$. Thus, either $n$ is a prime or it contains a new prime factor different from $p_i$, $1 \leq i \leq \ell$. This is a contradiction.
\end{proof}


\begin{algorithm}[H]
\caption{Sieve of Eratosthenes}
\label{alg:sieve}
\begin{algorithmic}[1]

\REQUIRE Integer $n \geq 0$
\ENSURE List of all prime numbers less than or equal to $n$

\IF{$n < 2$}
    \RETURN empty list $\varnothing$
\ENDIF

\STATE Initialize an array $\texttt{isPrime}[0\ldots n]$ to \TRUE
\STATE $\texttt{isPrime}[0] \gets \FALSE$
\STATE $\texttt{isPrime}[1] \gets \FALSE$

\FOR{$p \gets 2$ \TO $\lfloor\sqrt{n}\rfloor$}
    \IF{$\texttt{isPrime}[p] = \TRUE$}
        \FOR{$i \gets p^2$ \TO $n$ \textbf{step} $p$}
            \STATE $\texttt{isPrime}[i] \gets \FALSE$
        \ENDFOR
    \ENDIF
\ENDFOR

\STATE $\texttt{primes} \gets$ empty list $\varnothing$

\FOR{$i \gets 2$ \TO $n$}
    \IF{$\texttt{isPrime}[i] = \TRUE$}
        \STATE Append $i$ to $\texttt{primes}$
    \ENDIF
\ENDFOR

\RETURN $\texttt{primes}$

\end{algorithmic}
\end{algorithm}



Here is the python program for Seive of Erastothenes
\begin{minted}{python}
def sieve_of_eratosthenes(n: int) -> list[int]:
    """Generates a list of all prime numbers up to n."""
    if n < 2:
        return []

    # Initialize a tracking list with True values
    is_prime = [True] * (n + 1)
    is_prime[0] = is_prime[1] = False  # 0 and 1 are not prime

    # Only loop up to the square root of n
    for p in range(2, int(n**0.5) + 1):
        if is_prime[p]:
            # Mark all multiples of p starting from p^2 as False
            for i in range(p * p, n + 1, p):
                is_prime[i] = False

    # Extract all indices that remained True
    return [num for num, prime in enumerate(is_prime) if prime]


# Example Usage: Find all primes up to 30
print(sieve_of_eratosthenes(30))
# Output: [2, 3, 5, 7, 11, 13, 17, 19, 23, 29]
\end{minted}


\section{Prime Numbers}
The problem to decide whether a given number is a prime has been studied comprehensively. 
\begin{theorem}
    (The Prime Number Theorem) 
    \begin{align}
        \lim_{x \rightarrow \infty} \pi(x) \frac{\ln(x)}{x} = 1
    \end{align}
\end{theorem}


\begin{theorem}
    (Bertrand's postulate) For each positive integer $n > 1$ there is a
prime $p$ such that $n < p < 2n - 2$.
\end{theorem}

\begin{theorem}
    There are arbitrarily large gaps between consecutive primes.
\end{theorem}


\begin{proposition}
    The number of positive divisors of $n$ is
    \begin{align}
        d(n) &= (k_1 + 1)(k_2 + 1) \cdots (k_\ell +1)
    \end{align}
\end{proposition}


\begin{proposition}
    Let
    \begin{align}
        a &= p_1^{j_1} p_2^{j_2} \cdots p_{\ell}^{j_\ell} \\
        b &= p_1^{k_1} p_2^{k_2} \cdots p_{\ell}^{k_\ell} 
    \end{align}
    where $j_i, k_i \geq 0$ be two arbitrary positive integers. (We could assume that $a$ and $b$ are expressed using the same primes $p_1, p_2, \ldots , p_\ell$ because we allowed some exponents to be 0.) Then
    \begin{align}
        \gcd(a, b) &= p_1^{\min(j_1, k_1)} p_2^{\min(j_2, k_2)} \cdots p_\ell^{\min(j_\ell, k_\ell)} \\
        \lcm(a, b) &= p_1^{\max(j_1, k_1)} p_2^{\max(j_2, k_2)} \cdots p_\ell^{\max(j_\ell, k_\ell)} \\
    \end{align}
    Moreover,
    \begin{align}
        \gcd(a, b) \cdot \lcm(a, b) = a \cdot b
    \end{align}
\end{proposition}


\section{Modular arithmetic}

A nonempty set $I \subseteq \mathbb{Z}$ is called an \emph{ideal}\index{ideal} in $\mathbb{Z}$ if the following two conditions hold:
\begin{enumerate}
    \item For all $a, b \in I$, we have $a + b \in I$, and
    \item For all $k \in \mathbb{Z}$, and $a \in I$, we have $ka \in I$.
\end{enumerate}
In other words, the set $I$ is closed under addition, as well as under multiplication by any integer (and not just under multiplication by numbers belonging to $I$).

\begin{proposition}
    Every ideal in $\mathbb{Z}$ is principal.
\end{proposition}

A principal ideal $a\mathbb{Z}$ of $\mathbb{Z}$ will also sometimes be denoted by $\langle a\rangle$.

\begin{proposition}
    Let $a, b \in \mathbb{Z}$. Then $a \mathbb{Z} \subseteq b \mathbb{Z}$ if and only if $b|a$.
\end{proposition}


\begin{proposition}
    A number $a \in \mathbb{Z}/m\mathbb{Z}$ is a unit if and only if $\gcd(a,m)= 1$.
\end{proposition}

\section{Chinese Remainder Theorem}

\begin{theorem}
    Let $a$ and $b$ be (two) relatively prime numbers, $0 \leq r < a$ and $0 \leq s < b$. Then there exists a unique number $N$ such that $0 \leq N < ab$ and 
    \begin{align}
        r &\equiv N \pmod a \\
        s &\equiv N \pmod b
    \end{align}
    that is, $N$ has remainder $r$ on dividing by $a$ and remainder $s$ on dividing by $b$.
\end{theorem}


\section{Fermat's Little Theorem}


Fermat's little theorem is a special case of Euler's congruence theorem.
\begin{theorem}
    Let $p$ be a prime. If an integer $a$ is not divisible by $p$, then 
    \begin{align}
        a^{p-1} \equiv 1 \pmod p
    \end{align}
    In general, for all $a$, we have
    \begin{align}
        a^p \equiv a \pmod p
    \end{align} 
\end{theorem}
\begin{proof}
Page 89 of \cite{rubinstein2018cryptography}.
\end{proof}

\begin{example}
    The number $p = 15485863$ is prime, so Fermat's little theorem tells us that
    \begin{align}
        2^{15485862} \equiv 1 \pmod{15485863}.
    \end{align}
    Thus without doing any computing, we know that the number $2^{15485862}-1$ a number having more than two million digits, is a multiple of $15485863$.
\end{example}

\begin{definition}
   Let $n$ be an odd composite number and $b$ be a positive number such that $(b, n) = 1$. We call $n$, a \emph{Fermat pseudo-prime}\index{Fermat pseudo prime} under base $b$ iff 
   \begin{align}
    b^{n-1} \equiv 1 \pmod n
   \end{align} 
   Furthermore, we call $n$, \emph{Carmichael number} if it is Fermat pseudo-prime with base $b$ for all $b \in \mathbb{Z}_n^\times$.
\end{definition}


In 1922, the American mathematician, Robert Carmichael conjectured that
there are infinitely many Carmichael numbers. This was proved some seventy
years later by Alford, Granville, and Pomerance.

\begin{theorem}
    (Alford-Granville-Pomerance). The set of Carmichael numbers is infinite.
\end{theorem}
% W.R. Alford, A. Granville, and C. Pomerance, There are infinitely many Carmichael numbers, Annals of Mathematics 140 (1994), 703–722.



\section{Euler's Totient Function}
\begin{definition}
    Let $n$ be a positive integer. The number of positive integers, not exceeding $n$ and relatively prime to $n$, is denoted by $\phi(n)$. This function is called Euler's $\phi$-function or Euler's totient function.
\end{definition}

\begin{proposition}
    The Euler's totient function is multiplicative. That is, for $m, n \in \mathbb{Z}$ such that $\gcd(m, n) = 1$, we have 
    \begin{align}
        \phi(mn) = \phi(m) \phi(n)
    \end{align}
\end{proposition}

\begin{proposition}
    If $p$ is a prime number and $k$ is a positive integer, then 
    \begin{align}
        \phi(p^k) &= p^k - p^{k-1}
    \end{align}
\end{proposition}


\begin{proposition}
    Suppose $n$ factors as $n = \prod_{i=1}^\ell p_i^{k_i}$, then 
    \begin{align}
        \phi(n) &= n \prod_{i=1}^\ell \bigg( 1 - \frac{1}{p_i}\bigg)
    \end{align}
\end{proposition}


\section{Euler's congruence theorem}

\begin{theorem}
    (\emph{Euler's congruence theorem}) Let $n$ be a positive integer and $a$ be an integer relatively prime to $n$. Then
    \begin{align}
        a^{\phi(n)} \equiv 1 \pmod{n}
    \end{align}
\end{theorem}
\begin{proof}
    Page 90 of \cite{rubinstein2018cryptography}.
\end{proof}


\section{Legendre and Jacobi symbols}
Suppose $p$ is an odd prime and $a$ is an integer. The number $a$ is defined to be a \emph{quadratic residue modulo $p$} iff 
\begin{align}
    a \not\equiv 0 \pmod{p}
\end{align}
and 
\begin{align}
    x^2 \equiv a \pmod{p}
\end{align}
has a solution $x \in \mathbb{Z}_p$. The number $a$ is defined to be a \emph{quadratic non-residue modulo $p$} iff
\begin{align}
    a \not\equiv 0 \pmod{p}
\end{align}
and $a$ is not a quadratic residue modulo $p$.

\section{Primality testing}
% ref: Cryptography by Simon Rubinstein-Salzedo.pdf

RSA and other public key algorithms require large random prime numbers. The Prime
Number Theorem tells us that the asymptotic density of primes among the first $N$
integer numbers is $\dfrac{1}{\ln(N)}$. Fortunately, the density only decreases logarithmically.

\begin{example}
    The density is small, but not too small, since $\ln(x)$ increases slowly. For example, the density of primes among odd random numbers less than 22048 is approximately $\dfrac{2}{2048\ln(2)} \approx 0.0014$. The expected number of trials is $\dfrac{1}{0.0014} \approx 710$.
\end{example}

Given a number $N$ , how can we tell whether it is prime? One way of doing this is to divide it by all the numbers less than $N$ and see if it is exactly divisible by any of them. Or, we can be a little bit smarter: if $n$ factors as $N = ab$, then either $a$ or $b$ must be at most $\sqrt{N}$ , so we really only have to divide by numbers up to $\sqrt{N}$. This involves $O(\sqrt{N})$ steps. However, recall that $N$ only has $n = O(log(N ))$ digits, so $N = O(10^n)$. Hence, this algorithm for primality testing runs in $O(\sqrt{10^n})$ steps. Except not quite, because each ``step'' is a division, which itself requires something like $O(n^2)$ steps. Thus, overall, this algorithm for primality testing runs in $O(\sqrt{10^n} \times n^2)$ steps.

However, there are much faster algorithms for primality testing. The best provable algorithm is due to Agrawal, Kayal, and Saxena (\emph{AKS algorithm}), which in its
original form runs in $O(n^{12+\epsilon})$ steps, and more recent modifications have brought it down to $O(n^{6+\epsilon})$ steps. There are other algorithms that seem to do better, but that are hard to prove. For example, the Miller-Rabin algorithm, which is a probabilistic algorithm (i.e., one that relies on choosing random numbers and thus doesn't work every single time but that usually does) that can be turned into a deterministic algorithm if we assume that the Generalized Riemann Hypothesis is true, runs in time $O(n^{4+\epsilon})$ steps. The elliptic curve primality test seems to run in $O(n^{4+\epsilon})$ steps, but we are not able to prove this. 



% ref: RSA and Public Key Cryptography by Richard A.Molin

There do not exist any practical primality tests that encompass the three
desired properties of: 
\begin{enumerate}
    \item generality,
    \item speed, and 
    \item correctness.
\end{enumerate}


\subsection{Pepin's primality test}

{\bf Pepin's primality test}\index{Pepin's primality test}: Let $F_n = 2^{2^n} + 1$ be the $n$-th Fermat number. The Fermat's number $F_n$ is prime if and only if 
\begin{align}
3^{\frac{F_n - 1}{2}} \equiv - 1 \pmod{F_n}    
\end{align}


\subsection{Fermat Test}

In the 17th century, Pierre de Fermat proved the Fermat's Little Theorem, which can be turned into a simple primality testing algorithm.
Fermat's Little Theorem states that for any prime number $p$ and any number a not
divisible by $p$, the equivalence 
\begin{align}
    a^{p-1} \equiv \pmod{p}
\end{align}
must hold. Consequently, one can test the primality—or rather the compositeness—of $n$ by randomly choosing a value for $a$ (not divisible by $n$) and computing $a^{n-1} \pmod{n}$. If this value is not equal to 1, then $n$ is definitively not a prime (and we have found a witness for the compositeness of $n$, respectively). Unfortunately, the converse is not true and finding an $a$ for which $a^{n-1} \pmod{n}$ does not imply that $n$ is prime. In fact, there is an entire class of composite numbers for which $a^{n-1} \pmod{n}$ holds Pierre de Fermat was a French mathematician who lived from 1607 to 1665. For this reason, the Fermat test (and the two other tests mentioned later) is referred to as a compositeness test. Contemporary Cryptography for all a. These numbers are called Carmichael numbers. Because it is not able to correctly handle Carmichael numbers, the Fermat test is not widely deployed in practice. Instead, either the Solovay-Strassen test or the Miller-Rabin test is used.

\subsection{Solovay-Strassen test}

The Solovay-Strassen test is a probabilistic compositeness testing algorithm that was developed by Robert Solovay and Volker Strassen in 1976. It can prove the compositeness of a large number n with certainty, but it can prove the primality of $n$ only with a certain probability.

\subsection{Miller-Rabin test}


\begin{proposition}
    Let $n \in \mathbb{N}$ be odd and let $1 \leq a < n$ be an integer. If $\gcd(a, n) \neq 1$, then $n$ is composite. Otherwise, write
    \begin{align}
        n-1 &= 2^\ell m
    \end{align}
    with $k \in \mathbb{N}$ being maximal. If $n$ is prime, then exactly one of the following holds:
    \begin{enumerate}
        \item \begin{align}
        a^m \equiv \pm 1 \pmod{n}
    \end{align}
        \item there exists an integer $k \in \{1, \ldots, \ell-1\}$ such that
    \begin{align}
        a^{2^k m} \equiv -1 \pmod{n}
    \end{align}
    \end{enumerate}
\end{proposition}

If the Miller-Rabin algorithm outputs that a number is composite, this result must
be correct. However, there are bases $a$ such that a composite number $n$ is incorrectly identified as a prime in a run of the Miller-Rabin test.

\begin{algorithm}
\caption{Miller-Rabin Algorithm with $k$ runs}
\textbf{Input:} $n \in \mathbb{N}, k \in \mathbb{N}$ \\
\textbf{Output:} '$n$ is probably prime' or '$n$ is composite'
\begin{algorithmic}[1]
\FOR{$i = 1$ \TO $k$}
    \STATE $a \leftarrow^{\$} \{1, 2, \dots, n - 1\}$ \COMMENT{choose a uniform random base $a$.}
    \IF{$\gcd(a, n) \neq 1$}
        \STATE \textbf{return} '$n$ is composite'
    \ENDIF
    \STATE $y \equiv a^d \bmod n$
    \IF{$y \equiv 1$ \textbf{or} $y \equiv -1$}
        \STATE \textbf{do next} \COMMENT{$n$ could be a prime, test another base $a$.}
    \ENDIF
    \FOR{$r = 1$ \TO $s - 1$}
        \STATE $y = y^2 \bmod n$
        \IF{$y \equiv 1$}
            \STATE \textbf{return} '$n$ is composite' \COMMENT{$n$ cannot be a prime, stop.}
        \ENDIF
        \IF{$y \equiv -1$}
            \STATE \textbf{do next} \COMMENT{$n$ could be a prime, test another base $a$.}
        \ENDIF
    \ENDFOR
    \STATE \textbf{return} '$n$ is composite' \COMMENT{$n$ cannot be a prime, stop.}
\ENDFOR
\STATE \textbf{return} '$n$ is probably prime'
\end{algorithmic}
\end{algorithm}

\chapter{Polynomials}
\section{Ring of polynomials}
\section{Division algorithm}
\section{Lagrange's Interpolation}


\chapter{Groups}


\begin{proposition}
    Let $G$ is a finite group, $g \in G$, $k \in \mathbb{Z}$ is an integer, then
    \begin{align}
        o(g^k) &= \frac{o(g)}{(k, o(g))}
    \end{align}
    where the denominator is the greatest common divisor of $k$ and $o(g)$.
\end{proposition}
\begin{proof}
    Page 198 of \cite{zheng2022modern}.
\end{proof}

\begin{proposition}
    Suppose $G$ is a finite Abelian group, $a,b \in G$, and $(o(a),o(b))= 1$, then
    \begin{align}
        o(ab) &= o(a) o(b)
    \end{align}
\end{proposition}
\begin{proof}
    Page 198 of \cite{zheng2022modern}.
\end{proof}


\begin{theorem}
    (\emph{Existence of primitive root}) If and only if $n= 2, 4, p^k(k \geq 1)$ and $a = 2 p^k(k \geq 1)$ four cases, the primitive root of mod $n$ exists, where $p > 2$ is an odd prime.
\end{theorem}
\begin{proof}
    Page 199 of \cite{zheng2022modern}.
\end{proof}

\section{Permutation groups}

\section{Discrete logarithm}


\chapter{(Finite) Fields}

\begin{definition}
    A \emph{field}\index{field} $(\Bbbk, +, \cdot)$ is a set $\Bbbk$ with two binary operations
    \begin{align}
        \begin{array}{cccc}
        +: \Bbbk \times \Bbbk \longrightarrow \Bbbk & & &\cdot: \Bbbk \times \Bbbk \rightarrow \Bbbk  \\
         (r, s) \mapsto r+s & & & (r,s) \mapsto r\cdot s 
    \end{array}
    \end{align}
    called the \emph{field addition} and the \emph{field multiplication}, respectively, such that 
    \begin{enumerate}
        \item The pair $(\Bbbk, +)$ is an Abelian group, 
        \item The pair $(\Bbbk^\times, \cdot)$ is an Abelian group, and
        \item The field multiplication distributes over the field addition.
    \end{enumerate}
\end{definition}

\begin{example}
The field of rational numbers $\mathbb{Q}$.
\end{example}

\begin{example}
The field $\overline{\mathbb{Q}}$ of \emph{algebraic numbers}. Algebraic numbers  are complex numbers which are roots of some non-zero polynomial in one variable with rational coefficients. Therefore, by definition, it is the algebraic closure of $\mathbb{Q}$. 
\end{example}

\begin{example}
The field of real numbers $\mathbb{R}$. This can be constructed as the metric completion of the field of rational numbers with respect to the Euclidean metric.
\end{example}

\begin{example}
The field of complex numbers $\mathbb{C}$. It is the \emph{algebraic closure}\index{algebraic closure} of the real field $\mathbb{R}$.
\end{example}

\begin{example}
The field of $p$-adic numbers $\overline{\mathbb{Q}}_p$. This can be constructed as the metric completion of the field of rational numbers with respect to the metric induced by $p$-adic valuation given by $\nu_p(r) = k$ where $k$ is largest power of $p$ that divides the numerator of the reduced form of $r$. 
\end{example}

\begin{example}
The field of rational function $\mathbb{Q}(t)$. This consists of functions that are algebraic fractions such that both the numerator and the denominator are polynomials. 
\end{example}
    
\begin{example}
The set of all Laurent series with coefficients in $\mathbb{Q}$ is called the field of Laurent series over $\Bbbk$, and is denoted by $\Bbbk(\kern-.2em(t)\kern-.2em)$.
\end{example}

\begin{example}
If $\Bbbk$ is a field (such as $\mathbb{C}$), and $t$ an indeterminate, a \emph{Puiseux series}\index{Puiseux series} with coefficients in $\Bbbk$ is an expression of the form:
\begin{align}
    f = \sum_{n=m}^{\infty} c_n t^{ \frac{n}{k} }
\end{align}
where $k$ is a positive integer and $m$ is an integer. The set of all Puiseux series with coefficients in $\Bbbk$ is called the field of Puiseux series over $\Bbbk$, and is denoted by $\Bbbk\{\kern-.3em\{t\}\kern-.3em\}$. That is,
\begin{align}
    \Bbbk\{\kern-.3em\{t\}\kern-.3em\} &\defn \bigg\{ \sum_{n=m}^{\infty} c_n t^{ \frac{n}{k} }: ~\text{for all}~ k \in \mathbb{N}~\text{and for all}~ m \in \mathbb{Z} \bigg\}
\end{align}
It is the algebraic closure of field $\Bbbk(\kern-.2em(t)\kern-.2em)$ of formal Laurent series over $\Bbbk$.
\end{example}

\begin{theorem}
    Every finite integral domain is a field.
\end{theorem}


\begin{example}
    The ring $\mathbb{Z}_p = \{\overline{0}, \overline{1}, \cdots, \overline{p-1}\}$ for prime $p$ is a field with the usual operations of addition and multiplication modulo $p$.
\end{example}

A finite field of order $p^n$ for prime $p$ and positive integer $n$ is sometimes called a \emph{Galois field}, denoted $\GF(p^n)$ or $\mathbb{F}_{p^n}$.

\begin{proposition}
    If $p$ is a prime and $f(x)$ is an irreducible polynomial in $\mathbb{Z}_p[x]$ with $\deg(f) \geq 1$, then the quotient $\mathbb{F}_{p^n} \defn \dfrac{\mathbb{Z}_p[x]}{\langle f(x) \rangle}$ is a field with $p^n$, elements, where $n = \deg(f)$.
\end{proposition}

\begin{example}
    A field with nine elements: 
    \begin{align}
        \frac{\mathbb{Z}_3[x]}{x^2 + 1} \equiv \mathbb{Z}_3[i]
    \end{align}
\end{example}




% ref: Algebra for Cryptologists by Alko R.Meijer
\section{The multiplicative group}

\begin{proposition}
     The multiplicative group of a finite field is cyclic.
\end{proposition}

\chapter{Elliptic Curves}


\chapter{Probability}
% ref: Mathematics of Information by Stefan Schaffler
% ref: Probability by John W.Lamperti


\begin{definition}
    A collection of all elementary results, or outcomes of an experiment, is called a \emph{sample space}\index{sample space}. We usually denote it by $\Omega$.
\end{definition}


\innote{\caution}{In simple or finite settings, every single subset can be an event. In continuous or infinitely large sample spaces (like picking a random real number), assigning a probability to every single subset creates mathematical paradoxes.}


\begin{definition}
    A collection $\EuScript{S}$ of subsets of the sample space $\Omega$ is called a $\sigma$-algebra if all of the following are satisfied:
    \begin{enumerate}
        \item The sample space $\Omega \in \EuScript{S}$.
        \item If $S \in \EuScript{S}$, then $\Omega \backslash S \in \EuScript{S}$.
        \item If $S_1, S_2, \ldots,  \in \EuScript{S}$, then $\bigcup_{i=1}^\infty S_i \in \EuScript{S}$.
    \end{enumerate}
\end{definition}


A \emph{measurable space} $(\Omega, \EuScript{S})$ is a pair consisting of a \emph{sample space} $\Omega$ together with a $\sigma$-algebra $\EuScript{S}$ called the \emph{event space}. A \emph{probability space} $(\Omega, \EuScript{S}, \pr)$ is a triple consisting of a sample space $\Omega$, a $\sigma$-algebra $\EuScript{S}$ and a probability measure $\pr: \EuScript{S} \longrightarrow \mathbb{R}$ that satisfies:
\begin{enumerate}
    \item \emph{Non-negativity}: $\pr(S) \geq 0$, for all $S \in \EuScript{S}$.
    \item \emph{Normalization}: $\pr(\Omega) = 1$.
    \item \emph{Countable additivity}: Given disjoint sets $S_i \in \EuScript{S}$, we have 
    \begin{align}
        \pr\bigg( \bigcup_{i=0}^\infty S_i \bigg) = \sum_{i=0}^\infty \pr(S_i)
    \end{align}
\end{enumerate}

\chapter{Information Theory}
Information theory was founded by Claude Shannon in 1948 as a discipline to study the quantification, storage, and communication of information. In the past decades, it has played a critical role in modern communication as well as many other applications in engineering and computer science.  
% Mutual information and the Kullback-Leibler (KL) divergence 


\section{Information}
Let us consider the space of all messages to have a probability distribution. The probability distribution tell us how likely is it that we get a certain message. So we have some information about the occurrence of each message. Once we have receive a message, we gain some amount of information. For example, if probability of given message is 1, then we do not gain any new information after the occurrence of that message. The gain in the amount of information depends on its probability of its occurrence. Therefore, we can write the amount of information conveyed by a message $x$ as $I(p_x)$. This is also referred to as \emph{suprisal}, as it measures the amount of surprise we have after observing the message. \emph{Desidarata}: We would like the function $I: [0,1] \longrightarrow \mathbb{R}_{\geq 0}$ to satisfy the following properties:
\begin{enumerate}
    \item \emph{Non-negativity}: The information conveyed by a message cannot be negative. It has to be at least 0, that is
    \begin{align}
        I(p_x) \geq 0
    \end{align}
    % \item If the event is definite, that is, if the probability of the occurrence of the event is 1, then the corresponding information conveyed by the event will be 0. That is, 
    % \begin{align}
    %     I(1) &= 0
    % \end{align}
    \item \emph{Additivity}: The information conveyed by composite independent messages is simply given by the sum of the individual information content. That is,
    \begin{align}
        I(p_x p_y) &= I(p_x) + I(p_y)
    \end{align}
    \item \emph{Continuity}: The information $I$ is a continuous function.
\end{enumerate}


% ref: Python for Probability, Statistics and Machine Learning by Jose Unpingco
The only mathematical operator that satisfies these relations is the logarithm operator (up tp scalar multiple). Therefore, we have the following definition:
\begin{definition}
    The \emph{Shannon} \emph{information content} (or \emph{uncertainty content}) of an outcome $X = x$ is defined by
    \begin{align}
    I(x) &\defn \log\left(  \frac{1}{\pr(X = x)} \right) \\
    &= -\log(p(x))
    \end{align}
\end{definition}
The unit of the information is bits, nats or Hartley (or decits), when the base of the logarithm is the number $2$, the Euler's number $e$, or the number 10, respectively. The Shannon information can be interpreted as quantifying the level of \emph{surprise} of a particular outcome.



\section{Entropy}
% ref: The Probabilistic Method by Noga Alon and Joel H.Spencer
% ref: Statistical Properties of Deterministic Systems by Jiu Ding and Aihui Zhou
% ref: Permutation Complexity in Dynamical Systems by Jose M.Amigo
% ref: Pattern recognition and machine learning by Christopher M.Bishop
% ref: 

In cryptography, the security of most protocols and mechanisms depends on the generation of random sequences of bits or numbers. These sequences must have a sufficient length and be random in a very specific sense: we do not want an attacker to be able to guess part of or the whole sequence. To make this idea more precise, cryptographers use the concept of entropy.

Entropy is the amount of information provided by an observation of a random variable. It reflects the uncertainty associated with predicting the variable value. The larger the value of entropy, the greater the uncertainty in predicting the value of an observation. If $p_x = 1$ (the outcome is certain), then observing $x$ yields $-\log(1) = 0$ units of information. On the other hand, observing X = 0 in this case yields $-\log(0) = \infty$: total surprise at observing an impossible outcome.

\begin{definition}
    The (\emph{Shannon}) \emph{entropy}\index{entropy} or \emph{uncertainty} of a discrete random variable $X$ with a probability distribution function $\pr: \EuScript{S} \longrightarrow [0,1]$ is defined by
    \begin{align}
        H(X) &\defn \sum_{x \in \EuScript{S}} \pr(X = x) \log\left( \frac{1}{\pr(X = x)} \right)  \\
        \bigg[&= - \sum_{x \in \EuScript{S}} p_x \log(p_x) \bigg]\\
        &= \mathrm{E}[I(X)] 
    \end{align}
    If we restrict to binary logarithm $\log_2(\square)$, we refer to $H(X)$ as \emph{binary} (\emph{Shannon}) entropy. When the base is not mentioned, we assume that it to be 2. So mostly we are dealing with the binary entropy. We can also think of entropy as the average information. Observe that $\log(x)$ is undefined if $x = 0$. Hence, entropy is sometimes defined to be the relevant sum over all the non-zero probabilities\footnote{Also referred as $\supp(X)$}. However, since
        \begin{align}
        \lim_{x \longrightarrow 0} x \log\left( \frac{1}{x} \right) = 0
    \end{align}
    there is no real difficulty with allowing $\pr(x) = 0$ for some $x$'s. 
\end{definition}


\begin{example}
    Let the discrete random variable $X$ takes values from the set $\{1,2,3,4,5,6\}$ with $\pr(X = i) = \dfrac{1}{6}$ for each $i = 1, 2, \ldots, 6$. The entropy in this case is given by 
    \begin{align}
        H(X) &= -\sum_{i=1}^6 \frac{1}{6}\log(6) \\
            &= -\log(6)
    \end{align}
\end{example}

Let us look at the entropy of the various components of a cryptosystem. We
can think of the key as being a random variable $k$ that takes on values in $k$, and hence we can compute the entropy $H(k)$. Similarly, we can compute entropies $H(P)$ and $H(C)$ of random variables associated with the plaintext and ciphertext, respectively. 
\innote{}{\url{https://drand.love/}}
\innote{}{League of Entropy: \url{https://drand.love/loe/}}

\begin{example}
    Let $\EuScript{P} = \{0,1\}$ with 

\end{example}



\begin{example}
    If $X$ is binary with two possible values $x_1$ and $x_2$ such that
    \begin{align}
        \pr(X = x_1) &= p \\
        \pr(X = x_2) &= 1-p
    \end{align}
    then 
    \begin{align}
        H(X) = - p \log(p) - (1-p)\log(1-p)
    \end{align}
\end{example}


\begin{example}
    \emph{Entropy of English language}:
\end{example}


\begin{definition}
    A real-valued function $f$ is a concave function on an interval $I$ if
    \begin{align}
        \frac{f(x) + f(y)}{2} &\leq f \bigg( \frac{x+y}{2} \bigg)
    \end{align}
    for all $x, y \in I$. A function $f$ is a strictly concave function on an interval $I$ if
    \begin{align}
        \frac{f(x) + f(y)}{2} &< f \bigg( \frac{x+y}{2} \bigg)
    \end{align}
    for all $x, y \in I$, $x \neq y$.
\end{definition}

\begin{theorem}
    (\emph{Jensen's inequality}) Suppose $f$ is a continuous strictly concave function on the interval $I$. Suppose further that
    \begin{align}
        \sum_{i=1}^n a_i = 1
    \end{align}
    and $a_i > 0$, for $1 \leq i \leq n$. Then 
    \begin{align}
        \sum_{i=1}^n a_i f(x_i) &\leq f\bigg( \sum_{i=1}^n a_i x_i\bigg)
    \end{align}
    where $x_i \in I$, for $1 \leq i \leq n$. Further, equality occurs if and only if $x_1 = \ldots = x_n$.
\end{theorem}
% ref: Cryptography Theory and Practice by Douglas R.Stinson and Maura B.Paterson
We now proceed to derive several results on entropy

% There is also notion of entropy for (symbolic) dynamical systems for example logistic map.
Since entropy is a mathematical measure of uncertainty, it is natural to
expect that when $\pr(X =k) = \dfrac{1}{n}$, the corresponding value of $H$ is maximal.
\begin{proposition}
    Let $U_n$ denote the uniform distribution. 
    \begin{align}
        H(X) &\leq \log(n) 
    \end{align}
    for all $X$. The maximum is achieved if and only if $X \sim U_n$. In plain words, the maximum entropy is given if all the events occur equally often.
\end{proposition}
\begin{proof}
    
\end{proof}
% ref: An Introduction to Mathematical Cryptography by Jeffrey Hoffstein, Jill Pipher and Joseph H.Silverman
% ref: Basic Principles and Applications of Probability Theory by A.V.Skorokhod

\begin{proposition}
If a function $f: \EuScript{S}_n \longrightarrow [0, \infty]$ satisfy the requirements of continuity, weak additivity, monotonicity, branching and normalization, then
\begin{align}
    f(X) &= H(X)
\end{align}
\end{proposition}

Application of these include \emph{Huffman coding} and \emph{Error correcting codes}. Huffman codes are used in the JPEG and MPEG standards.
% ref: Cryptography and Data Security by Dorothy Elizabeth Robling Denning


\begin{definition}
    Let $X$ and $Y$ be two discrete random variables with the event space $\EuScript{S}$ and $\EuScript{T}$, respectively. The \emph{joint entropy} $H(X, Y)$ is defined as
    \begin{align}
        H(X, Y) &\defn \sum_{ \substack{x \in \EuScript{S} \\ y \in \EuScript{T}}} \pr(X = x, Y = y) \log\bigg( \frac{1}{\pr(X = x, Y = y)}\bigg) \\
        &= \sum_{ \substack{x \in \EuScript{S} \\ y \in \EuScript{T}}} p(x, y) \log \bigg( \frac{1}{p(x, y)}\bigg)
    \end{align}
    The \emph{conditional entropy} $H(X|Y)$ is defined as 
    \begin{align}
        H(X|Y) &\defn \sum_{y \in \EuScript{T}} p(y)H(X|Y=y) \\
        &= \sum_{\substack{x \in \EuScript{S} \\ y \in \EuScript{T}}} p(x, y) \log \bigg( \frac{1}{p(x|y)}\bigg)
    \end{align}
\end{definition}


% ref: Cryptography Theory and Practice by Douglas R.Stinson and Maura B.Paterson
\begin{theorem}
    Let $\mathscr{C} = (\EuScript{P}, \EuScript{C}, \EuScript{K}, \EuScript{E}, \EuScript{D})$ be a cryptosystem, then 
    \begin{align}
        H(K|C) = H(K) + H(P) - H(C)
    \end{align}
\end{theorem}


% ref: An Introduction to Mathematical Cryptography
\begin{definition}
    Let $L$ be a language and for each $n \geq 0$, let $L^n$ denote\footnote{Abuse of notation.} the random variables whose values are strings of $n$ consecutive characters of $\mathsf{L}$. The entropy of L is defined to be the quantity
\begin{align}
    H(L) &\defn \lim_{n \rightarrow \infty} \frac{H(L^n)}{n}
\end{align}
\cite{silverman2008introduction}
\end{definition}

Although it is not possible to precisely determine the entropy of the English language $E$, experimentally it appears that
\begin{align}
    1.0 \leq H(E) \leq 1.5
\end{align}
This means that despite the fact that it requires almost five bits to represent each of the 26 letters used in English, each letter conveys less than one and a half bits of information. Thus English is approximately 70\% redundant\footnote{This does not mean that one can remove 70\% of the letters and still have an intelligible message. What it means is that in principle, it is possible to take a long message that requires 4.7 bits to specify each letter and to compress it into a form that takes only 30\% as many bits.}.
\section{Kullback-Leibler divergence}



\section{Information entropy analysis}
The entropy analysis is used to test the robustness of the encryption algorithm. It is very important to compare the entropy of plain images and ciphered images. By doing so, we are able to see that the entropy of the encrypted images is approximately equal to an 8-bit depth. This will help us to prove the encryption technique against the entropy attack

\section{Perfect Secrecy}
A cryptosystem has perfect secrecy if the interception of a ciphertext gives
the cryptanalyst no information about the underlying plaintext and no in-
formation about any future encrypted messages.



% ref: Cryptography Theory and Practice by Douglas R.Stinson and Maura B.Paterson
\begin{definition}
    A cryptosystem has perfect secrecy if 
    \begin{align}
        \pr(x|y) = \pr(x)
    \end{align}
    for all $x \in \EuScript{P}$, $y \in \EuScript{C}$. That is, the \emph{a posteriori} probability that the plaintext is $x$, given that the ciphertext $y$ is observed, is identical to the \emph{a priori} probability that the plaintext is $x$.
\end{definition}








% \chapter{Cryptographic  Fundamentals}

% \emph{Cryptography} is the practice and study of techniques for secure communication in the presence of adverse third parties. \emph{Encryption} refers to the process of changing ordinary text (plaintext) into unintelligible text (ciphertext). \emph{Decryption} is the reverse.

% references: 
% 1. Introduction to Cryptography with Maple by Jose Luis Gomez Pardo
% 2. Introduction to Cybersecurity by Robin Sharp
% 3. Understanding Cryptography by Christof Paar, Jan Pelzl and Tim Guneysu
% 4. Introduction to Cryptography by Hans Delfs and Helmut Knebl
% 5. Introduction to Computer Networks and Cybersecurity by Chwan-Hwa Wu and J.David Irwin
% 6. Handbook of Financial Cryptography and Security by Biurton Rosenberg
% 7. Guide to Internet Cryptography by Jorg Schwenk
% 8. Guide to Computer Network Security by Joseph Migga Kizza
% 9. Encrypt, Sign, Attack by Olaf Manz
% 10. Cryptograpy Theory and Practice by Douglas R.Stinson
% 11. Discovering Cybersecurity by Seth James Nielson
% 12. CryptoSchool by Joachim von zur Gathen copy
% 13. Cryptology by Richard Klima and Neil Sigmon
% 14. Cryptography by Douglas R.Stinson and Maura B.Paterson
% 15. Cryptography and Network Security by William Stallings 
% 16. Cryptographic Protocol by Ling Dong and Kefei Chen
% 17. Cryptography and Network Security by Marcelo Sampaio de Alencar
% 18. Computer Security Fundamentals by Chuck Easttom
% 19. Computer Security and the Internet by Paul C.van Oorschot
% 20. Computer Security and Cryptography by Alan G.Konheim.pdf
% 21. CompTIA Security Certification Study Guide by Ahmed F.Sheikh
% 22. Classical and Modern Cryptography for Beginners by Rajkumar Banoth and Rekha Regar
% 23. A Classical Introduction to Cryptography by Serge Vaudenay













































\chapter{Hash Function}

These are one-way function\footnote{How can we make it rigorous.}. We desire the following properties from these functions, (call an arbitrary one $f$):
\begin{enumerate}
    \item It can be applied to a block of data of any size.
    \item It produces a fixed-length output.
    \item The value $f(x)$ is relatively easy to compute for any given $x$.
    \item \emph{One-way}: For any given value $h$, it is computationaly infeasible to find (any one of) its pre-image block $x$, that is such that $f(x) = h$.
    \item \emph{Weak collision resistant}: For any given block $x$, it is computationally infeasible to find $x' \neq x$ such that $f(x) = f(x')$.
    \item \emph{Strong collision resistant}: It is computationally infeasible to find any pair $(x, x')$ such that $f(x) = f(x')$.
\end{enumerate} 


Useful applications of hash function:
\begin{enumerate}
    \item Human-readable method to compare/verify data.
    \item Chaining events together - blockchain.
    \item Digital signatures and message authentication codes.
    \item Fundamental building block of many secure protocols.
\end{enumerate}

\begin{example}
    \emph{SHA-0}: Completely broken, collisions found.
    \emph{SHA-1}: Implemented using Merkle-Damg\.ard construction. 
    \emph{SHA-2}: No known flaws.
    \emph{SHA-3}: Subset of Keccak. Implemented using sponge construction.
\end{example}



\begin{example}
    \emph{MD5}: Completely broken, collisions found.
\end{example}

\chapter{Codes}

We used the reference \cite{zheng2022modern} for this chapter. Let $n$ be a positive integer, $\mathbb{F}_q^n$ is an $n$-dimensional linear space over $\mathbb{F}_q$ , also called \emph{codeword space}\index{codeword space}. A vector $\mathbf{x} = (x_1, x_2, \ldots , x_n)$ in $\mathbb{F}_q^n$ is called a \emph{codeword of length $n$}. For convenience, a codeword $\mathbf{x}$, we write as $\mathbf{x} = x_1x_2\cdots x_n$, each $x_i \in \mathbb{F}_q$ is called a \emph{character}. We denote by $0 = (0,0, \ldots, 0)$.
Given codewords $\mathbf{x} = x_1 x_2 \cdots x_n$ and $\mathbf{y} = y_1 y_2 \cdots y_n$, we define the \emph{Hamming distance}\index{Hamming distance} $d(\mathbf{x}, \mathbf{y})$ by 
\begin{align}
    d(\mathbf{x}, \mathbf{y}) &\defn \# \bigg\{ i: 1 \leq i \leq n~\st~x_i \neq y_i\bigg\}
\end{align}
Obviously
\begin{align}
    0 \leq d(\mathbf{x}, \mathbf{y}) \leq n
\end{align}
The \emph{weight function}\index{weight function} $\wt: \mathbb{F}_q^n \longrightarrow \mathbb{N}$ is defined by 
\begin{align}
    \wt(\mathbf{x}) &\defn d(x, 0)
\end{align}
that is, the Hamming distance between $\mathbf{x}$ and $0$. The following properties are obvious:

\begin{proposition}
    If $\mathbf{x}, \mathbf{y} \in \mathbb{F}_q^n$, then 
    \begin{enumerate}
        \item \emph{Non-negativity}: $d(\mathbf{x}, \mathbf{y}) \geq 0$, where equality holds if and only if $\mathbf{x} = \mathbf{y}$.
        \item \emph{Symmetry}: $d(\mathbf{x}, \mathbf{y}) = d(\mathbf{y}, \mathbf{x})$.
        \item $\wt(-\mathbf{x}) = \wt(\mathbf{x})$.
        \item \emph{Translation invariance}: $d(\mathbf{x}, \mathbf{y}) = d(\mathbf{x} - \mathbf{z}, \mathbf{y} - \mathbf{z})$ for all $\mathbf{z} \in \mathbb{F}_q^n$.
        \item $d(\mathbf{x}, \mathbf{y}) = \wt(\mathbf{x} - \mathbf{y})$.
    \end{enumerate}
\end{proposition}

\begin{proposition}
    Let $\mathbf{x}, \mathbf{y} \in \mathbb{F}_q^n$ be two codewords, then
    \begin{align}
        \wt(\mathbf{x} \pm \mathbf{y}) \leq \wt(\mathbf{x}) + \wt(\mathbf{y})
    \end{align}
\end{proposition}

\begin{definition}
    A \emph{linear code}\index{linear code} $C$ is a vector subspace of $\mathbb{F}_q^n$.
\end{definition}


\section{Error Correcting Codes}
\section{Hamming Codes}
\section{Polynomial Codes}
\section{Reed-Solomon (RS) Codes}


\chapter{Complexity}
Big Oh, Big Omega, Big Theta













% \chapter{References}

%-- Computation Engineering by Ganesh Gopalakrishnan

% \cite{adamson2012set}

% \nocite{*}


\bibliography{../Bibliography/References}
\bibliographystyle{amsalpha}

\printindex




\end{document}